\section{Substitutions and Almost Symmetric Functions}

Substition refers to the set of permutations of $n$ elements which form the Symmetric Group.

\begin{theorem}
All the substitutions of $x_1, x_2, \hdots, x_n$ which leave unaltered a rational function
$\phi(x_1, x_2,\hdots, x_n$ form a group $G$.
\end{theorem}

Let $\phi_s$ denote the function obtained by applying to $\phi$ the substitusion $s$. If $a$ and $b$
are two substitutions which leave $\phi$ unaltered, then $\phi_a = \phi, \phi_b = \phi$. Then

\[ (\phi_a)_b = (\phi)_b = \phi_b = \phi, \mbox{ or } \phi_{ab} = \phi \]
Hence the product of $ab$ is one of the substitutions which leave $\phi$ unaltered and the set
$G$ has the group property. \\

The group $G$ is called the {\it group of the function $\phi$}, while $\phi$ is said to {\it belong to the
group $G$}\\

\begin{theorem}
Every substitution can be expressed as a product of transpositions in various ways.
\end{theorem}

\begin{theorem}
Of the various decompositions of a given substitution $s$ into a product of transpositions, all contain
an even number of transpositions (an even substitution), or all constain an odd number of transpositions
(an odd substitution).
\end{theorem}

\begin{corollary}
The totality of even substitutions on $n$ letters forms a group, called the Alternating Group on n letters.
\end{corollary}


Consider $n$ elements $x=\{ x_1, x_2, x_3, \hdots, x_n \}$. Consider the rational function

$$\displaystyle
\begin{array}{r}
\phi(x) = \prod_{i < j}\, (x_i - x_j) = (x_1-x_2) (x_1-x_3) (x_1-x_4) \quad \hdots\hdots\quad (x_1 - x_n) \\
(x_2-x_3) (x_2-x_4) \quad \hdots\hdots\quad (x_2 - x_n) \\
(x_3-x_4) \quad \hdots\hdots\quad (x_3 - x_n) \\
\quad\quad \hdots\hdots\hdots \quad\quad\quad~~\\
(x_{n-1} - x_n) \\
\end{array}
$$
consisting of the product of all the difference of the elements.\\

$\phi(x)$ is related to the Vandemonde determinant

$$\phi(x) = V(x) =
\begin{vmatrix}
1 & x_1 & x_1^2 & x_1^3 & \hdots & x_1^{n-1} \\
1 & x_2 & x_2^2 & x_2^3 &\hdots &x_2^{n-1} \\
1 & x_3 & x_3^2 & x_3^3 &\hdots & x_2^{n-1} \\
\vdots \\
1 & x_n & x_n^2 & x_n^3 &\hdots & x_n^{n-1} \\
\end{vmatrix}
$$

The square of $V(x)$ is a
well- known symmetric function, the discriminant $\triangle$. Therefore $V(x)$ has two values equal
numerically with opposites signs, viz. $\sqrt{\triangle}$ and $-\sqrt{\triangle}$. $V(x)$ is an example
of an alternating function. Any transposition alters the sign of $V(x)$. For example, the transposition
$(x_1x_2)$ affects only the terms in the first and second lines of the product, and replaces them by
\[
\begin{array}{r}
(x_2-x_1) (x_2 - x_3) (x_2 - x_4) \quad \hdots\hdots\quad (x_2 - x_n) \\
(x_1-x_3) (x_1-x_4) \quad \hdots\hdots\quad (x_1 - x_n). \\
\end{array}
\]
Hence, if $s$ is the product of an even number of transpositions, it leaves $v(x)$ unaltered; if $s$
is the product of an odd number of transpositions, it changes $V(x)$ into $-V(x)$. \\